\documentclass[11pt]{article}
\usepackage[a4paper,margin=22mm]{geometry}
\usepackage{fontspec}
\setmainfont{DejaVu Serif}
\usepackage{amsmath,amssymb,amsthm,mathtools,graphicx,xcolor,enumitem,hyperref,xurl}
\setlength{\emergencystretch}{4em}
\newtheorem{theo}{Theorem}[section]
\newtheorem{lemm}[theo]{Lemma}
\newtheorem{prop}[theo]{Proposition}
\newtheorem{coro}[theo]{Corollary}
\newtheorem{corol}[theo]{Corollary}
\theoremstyle{definition}
\newtheorem{defi}[theo]{Definition}
\newtheorem{exam}[theo]{Example}
\newtheorem{exem}[theo]{Example}
\newtheorem{rema}[theo]{Remark}
\newtheorem*{rema*}{Remark}
\newtheorem*{defi*}{Definition}
\newtheorem*{theo*}{Theorem}
\newtheorem*{lemm*}{Lemma}
\newtheorem*{prop*}{Proposition}
\newcommand{\enoncename}{}
\newtheorem{corpusenonce}[theo]{\enoncename}
\newenvironment{enonce}[1]{\renewcommand{\enoncename}{#1}\begin{corpusenonce}}{\end{corpusenonce}}
\newenvironment{enonce*}[1]{\par\medskip\noindent\textbf{#1.}\itshape}{\par\medskip}
\newenvironment{proofc}{\begin{proof}}{\end{proof}}
\newcommand{\Mc}{\mathcal{M}}

\numberwithin{equation}{section}
%~Mouliné par MaN_auto v.0.21.8 2019-10-30 13:57:13


\usepackage{mathrsfs}

\DeclareMathOperator{\Isom}{Isom}
\DeclareMathOperator{\Aff}{Aff}
\DeclareMathOperator{\Co}{Cone}
\DeclareMathOperator{\Cay}{Cay}
\DeclareMathOperator{\Proba}{Proba}
\DeclareMathOperator{\supp}{supp}








\newcommand{\A}{\measuredangle}
\newcommand*{\bfZ}{\mathbf{Z}}

\newenvironment*{step}{\begin{proof}}{\let\qed\relax\end{proof}}
\newenvironment*{laststep}{\begin{proof}}{\end{proof}}

\newtheorem{mthm}{Theorem}
\newtheorem{mcoro}[mthm]{Corollary}
\renewcommand*{\themthm}{\Alph{mthm}}

\graphicspath{{./figures/}}

\newcommand*{\mk}{\mkern -1mu}
\newcommand*{\Mk}{\mkern -2mu}
\newcommand*{\mK}{\mkern 1mu}
\newcommand*{\MK}{\mkern 2mu}

%\hypersetup{urlcolor=purple, linkcolor=blue, citecolor=red}

\newcommand*{\romanenumi}{\renewcommand*{\theenumi}{\roman{enumi}}}
\newcommand*{\Romanenumi}{\renewcommand*{\theenumi}{\Roman{enumi}}}
\newcommand*{\alphenumi}{\renewcommand*{\theenumi}{\alph{enumi}}}
\newcommand*{\Alphenumi}{\renewcommand*{\theenumi}{\Alph{enumi}}}


\begin{document}
\section*{2594: Relatively proper affine isometric cocycles in mixed lp–l1 normed spaces for finitely generated relatively hyperbolic groups}
\noindent\textbf{Author:} Indira Chatterji; François Dahmani\par\noindent\textbf{Source:} Proper actions on \$\textbackslash{}ell \textasciicircum{}p\$-spaces for relatively hyperbolic groups. Annales Henri Lebesgue 3 (2020), publisher version, DOI 10.5802/ahl.26. \url{https://ahl.centre-mersenne.org/item/AHL_2020__3__35_0/}
\par\noindent\textbf{License:} CC BY 4.0, \url{https://creativecommons.org/licenses/by/4.0/}. The original publisher PDF has the explicit license notice. See the saved source/license records.
\par\noindent\textbf{Editorial material note (not author text):} The selected problem is Theorem 3.2 in the authored mixed lp--l1 normed space. Complete author Sections 1--3 retain all new conical-geodesic flow, confluence, mask and summability support. The later subgroup properness and stronger pure-lp claims are not selected. All proof text and formulas below are editable author TeX. Mathematical wording, language and proof order are retained. Only the article class, title matter and bibliography presentation are adapted for independent compilation; no proof has been generated or repaired.
\par\noindent\textbf{Editorial difficulty:} H2: The authors construct a flow on nonlocally finite fine hyperbolic graphs, obtain finite conical slices and checkpoint collapse at large angles, prove exponential confluence against word growth and disjoint-mask lower bounds, then assemble an equivariant summable cocycle. This new geometric-group construction is beyond ordinary qualifying examinations.
\medskip\hrule\medskip
\noindent\textbf{Author statement, complete proof and local supporting text follow in source order.}
\section{Working in relatively hyperbolic groups}\label{bases}

To fix notations we recall that a \emph{graph} $X$ is a set $X^{(0)}$, the \emph{vertex set}, with a set $X^{(1)}\subseteq X^{(0)}\times X^{(0)}$, the \emph{oriented edges}, endowed with an \emph{origin map} $o:X^{(1)} \to X^{(0)}$ and a fixed-point free involution reversing the edges $\bar{\ }: X^{(1)}\to X^{(1)}$. A pair $\{e,\bar e\}$ is an \emph{unoriented edge}. We denote by $t:X^{(1)} \to X^{(0)}$ the composition $\bar o$, this is the \emph{terminus map}. We assume the reader familiar with paths, connectedness, length, geodesics, in graphs, as well with hyperbolicity. All our graphs are considered with their graph metric where an edge has length one.

\begin{defi}
Let $X$ be a graph. Given a vertex $v$ and two oriented edges $e_1, e_2$ such that $o(e_2) = t(e_1) =v$, the \emph{angle between $e_1$ and $e_2$ at $v$}, denoted by $\A_v(e_1,e_2)$, is the infimum (in $\mathbb{R_+} \cup \{+\infty\}$) of length of paths from $o(e_1)$ to $t(e_2) $ that avoid $v$. In other words,
\[
\A_v(e_1,e_2)=d_{X\setminus\{v\}}(o(e_1),t(e_2)).
\]
By convention, we will use the following abuse of notation: if $x_1, x_2,c$ are vertices in $X$, we say that $\A_c(x_1,x_2)>\theta$ if there exists two edges $e_1, e_2$ with $t(e_1)= o(e_2) =c$, with $e_i$ on a geodesic between $c$ and $x_i$ ($i=1,2$) and such that $\A_c(e_1,e_2)>\theta$. We say that $\A_c(x_1,x_2)\leq\theta$ otherwise, namely, if for any geodesic between $c$ and $x_i$ and starting by and edge $e_i$($i=1,2$), then $\A_c(e_1,e_2)\leq\theta$.
\end{defi}

Large angles at a vertex in a $\delta$-hyperbolic graph will force geodesics through that vertex. More precisely we have the following.

\begin{prop}\label{prop;goulet}
Let $X$ be a $\delta$-hyperbolic graph, and let $a,b,c\in X^{(0)}$.
\begin{enumerate}
\item \label{1.2.1} If $\A_c{(a,b)} >12\delta$, then all geodesics from $a$ to $b$ contain $c$.
\item \label{1.2.2} If $e,e'$ are two edges originating at $c$ and that are on geodesics from $c$ to $a$, then $\A_c(e, \bar e') \leq 6\delta$.
\item \label{1.2.3} For any $\theta\geq 0$, if $\A_c(x,y) > \theta + 12\delta$ then for all choice of edges $\epsilon_1, \epsilon_2$ at $c$ starting geodesics respectively to $x$ and to $y$, $\A_c(\epsilon_1, \bar \epsilon_2) \geq \theta$.
\end{enumerate}
\end{prop}

\begin{proof}
\begin{step}[\eqref{1.2.1}]
Applying $\delta$-thinness of triangles on the
geodesics $[c,a]$ and $[c, b]$ at distance $2\delta$ from $c$, we obtain a path of length at most $12\delta$, containing an arc of $[a,b]$, and that goes from the end of the first edge of $[c,a]$ to the end of the first edge of $[c,b]$. If the angle at $c$ is larger than $12\delta$ this path cannot avoid $c$. But only the arc on $[a,b]$ can meet $c$ again. Hence the first claim.
\end{step}

\begin{step}[\eqref{1.2.2}]
Any geodesic bigon between $c$ and $a$ must be $\delta$-thin, so if we take a geodesic bigon starting with the edges $e$ and $e'$ and look at the points at distance $2\delta$ from $c$, we get a path of length $6\delta$ at most and that avoids $c$.
\end{step}

\begin{laststep}[\eqref{1.2.3}]
This is a consequence of~\eqref{1.2.1} and~\eqref{1.2.2}, combined with the triangular inequality for angles.
\end{laststep}
\let\qed\relax
\end{proof}

\begin{defi}
A graph is called \emph{fine} if for each edge $e$, and for all $L\geq 0$, the set of edges
\[
\{ e'\,|\,t(e')=o(e), \A_{o(e)}(e, e') \leq L\}
\]
is finite.
\end{defi}

This definition of fine graphs is equivalent to Bowditch's definition that, for all edge, for all number, there are only finitely many simple loops of this length passing through this edge, \cite{BowRH}.


\begin{defi}
Let $G$ be a finitely generated group, and $H_1, \dots, H_k$,
subgroups of $G$. Consider a Cayley graph $\Cay(G)$ of $G$ (over a finite generating set), and construct the coned-off over $H_1, \dots, H_k$ as follows: for all $i$ and all left coset $gH_i$ of $H_i$, add a vertex $\widehat{gH_i}$ and for each $h\in H_i$, add an edge between $\widehat{gH_i}$ and $gh$. We denote by $\widehat{{\Cay}} (G)$ this graph, called a \emph{coned-off Cayley graph} (of $G$ with respect to the subgroups $H_1, \dots, H_k$).
\end{defi}

We can now recall the definition of Bowditch of relative hyperbolicity from~\cite{BowRH} (his original definition is about $G$-graphs, but this is an equivalent formulation from~\cite[Appendix]{Dthesis}).

\begin{defi}
A pair $(G,\{H_1, \dots, H_k\})$ of a group $G$ with a collection of subgroups, is \emph{relatively hyperbolic}, if a (equivalently any) coned-off Cayley graph $\widehat{{\Cay}} (G)$ over $H_1, \dots,H_k$ is hyperbolic and fine. The groups $H_i$ are called \emph{peripheral subgroups}.
\end{defi}

A difficulty of working in relatively hyperbolic groups is the lack of local finiteness in the coned-off graph. The angles, and below the cones, are useful tools for working around this. The following estimate says that the word metric is equivalent to the coned-off metric to which we add the sum of the angles at infinite valence vertices. We will be needing the first inequality only.

\begin{prop}\label{prop;pseudo_word_metric}
Let $G$ be a group hyperbolic relative to $\{H_1, \dots, H_k\}$ with $d_w$ a word metric, for a finite generating set containing generating sets of each $H_i$. Let $\widehat{{\Cay}} (G)$ be its coned-off Cayley graph, for its relatively hyperbolic structure, with $\hat{d}$ the graph metric.

There are $A,B>1$ for which, for all $g\in G$, for each choice of geodesic $[1,g]$ in the coned-off graph $\widehat{{\Cay}} (G)$, if $\sum\A ([1,g])$ denotes the sum of angles of edges of $[1,g]$ at vertices of infinite valence, then
\[
\frac1A d_w(1,g) -B \leq \hat{d}(1,g) + \sum\A ([1,g]) \leq A\,d_w(1,g) +B.
\]
\end{prop}

\begin{proof}
First, given a geodesic $[1,g]$ in the coned-off graph $\widehat{{\Cay}} (G)$, we construct a path in the Cayley graph (whose length will bounded from above $d_w(1,g)$) by replacing each pair of consecutive edges $e,e'$ at an infinite valence vertex, by a path in the corresponding coset $xH_i$ of $H_i$. We remark that we can choose this path of length bounded above in terms of the angle, say $\theta$. Indeed, by definition of angle, there exists a path in $\widehat{{\Cay}} (G)$ of length at most $\theta$ avoiding $\widehat{xH_i}$, from $o(e)$ to $t(e')$. Project each vertex in this path on the $1$-neighborhood of $\widehat{xH_i}$. One gets a sequence of $\theta$ points $y_1, \dots y_\theta$ in $xH_i$, from $o(e)$ to $t(e')$, two consecutive points being at an angle of at most $2\delta$ at $\widehat{xH_i}$. It follows that there are only boundedly many possible transitions $y_i^{-1}y_{i+1}$ in $H_i$, and we have our bound on the word length in $H_i$ in term of the angle $\theta$ at $\widehat{xH_i}$. Thus $ \frac1A d_w(1,g) -B \leq \hat{d}(1,g) + \sum\A ([1,g]) $ (for some $A$ and~$B$).

For the second inequality, consider a geodesic from $1$ to $g$ in the word metric. It produces an injective continuous path $q$ in $\widehat{{\Cay}} (G)$, that fellow-travels a geodesic $[1,g]$ in the coned-off graph (see for instance~\cite[Prop.~8.25]{DS}). The path $q$ does not contain any vertex of infinite valence, therefore each time $[1,g]$ contains a vertex of infinite valence, we can use the proximity of $[1,g]$ with $q$ and the fact that $q$ does not contain this point, to estimate the angles in terms of the length of a subsegment of $q$. By thinness of the graph, each subsegment of $q$ is only used at most a uniformly bounded amount of times. Therefore, the total sum of the angles at infinite valence points is at most a certain multiple of the length of $q$, which is the word length of~$g$.
\end{proof}

\subsection{Cones}
The notion of angles will now allow us to talk about cones in any graph $X$. Cones are a useful tool for working in relatively hyperbolic groups as they allow to take neighborhoods of geodesics in a finite way.

\begin{defi}
Let $X$ be a graph. Given an oriented edge $e$ and a number $\theta>0$, the \emph{cone of parameter $\theta$ around the oriented edge $e$}, denoted $\Co_{\theta}(e)$ is the subset of vertices $v$ and edges $\epsilon$ of $X$ such that there is a path from $o(e)$, starting by $e$, and containing $v$ or $\epsilon$, that is of length at most $\theta$ and for which two consecutive edges make an angle at most $\theta$.

For vertices of finite valence we can define the cone at a vertex of parameter $\theta$ by the union of all the cones of parameter $\theta$ around the edges adjacent to the vertex.
\end{defi}

If angles inside a cone of parameter $\theta$ are by definition bounded by $\theta$, angles at vertices close to the cone are also controlled in terms of $\theta$. Precisely we have the following estimate.

\begin{lemm}\label{lem;thetasquare}
In any graph $X$, for any oriented edge $e$ and $\theta\geq 0$, if $v\in \Co_\theta(e)$ then, for all $c\in X^{(0)}$ at distance $\leq
\theta/10$ from both $v$ and $e$, then $\A_c(o(e), v)\leq (\theta^2+3\theta)/2$.
\end{lemm}

\begin{proof}
Consider $c$ with two edges $\epsilon_1,
\epsilon_2$ starting at $c$, on geodesics $g_1, g_2$ toward $o(e)$, and $v$ respectively, of length $\leq \theta/10$. We also are given a path $p$ starting by $e$, going to $v$ of length at most $\theta$, and maximal angle at most $\theta$ between consecutive edges. Concatenating the paths $g_1$, $p$, $\bar g_2$ (orientation reversed), we have a loop at $c$, starting by $\epsilon_1$, and ending by $\epsilon_2$ and of length at most $6\theta/5$. If this loop doesn't pass by $c$ (except starting and ending point), then by definition of angle we have that $\A_c(
\epsilon_1, \bar \epsilon_2) \leq 6\theta/5$, which would prove the claim. If this path passes at least twice at $c$, only the segment $p$ can use $c$, and it can do so at most $(\theta- d(o(e),c)- d(c,v))/{2}$ times, each time realizing an angle between incident and exiting edge of at most $\theta$. By the triangular inequality on angles, this path gives a bound on the angle $\A_c(\epsilon_1, \bar \epsilon_2)$ of at most $\theta(\theta- d(o(e),c)- d(c,v))/{2}+ 6\theta/5\leq (\theta^2+3\theta)/2$. This establishes the claim.
\end{proof}

Cones also can be composed in an obvious way.

\begin{prop}\label{prop;composition_of_cones}
Let $\alpha,\beta\geq 0$, if $e,e', e''$ are three edges of a graph and if $e''\in
\Co_\alpha(e')$ and $e'\in \Co_\beta(e)$ then $e''\in \Co_{\alpha+\beta}(e)$.
\end{prop}
\begin{proof}
We have a path of length and maximal angle bounded by $\beta$ that starts by $e$ and ends by $e'$ or $\bar e'$, and another of length and maximal angle at most $\beta$ that starts at $e'$ and contains $e''$ or $\bar e''$. Concatenate the two paths, possibly overlapping at $e'$, or possibly cancelling a backtrack $e'\bar e'$, produces a path of length at most $\alpha +\beta$ and maximal angle at most $\max \{\alpha, \beta\}$.
\end{proof}

Finally, the following proposition shows that cones are well-suited to work in hyperbolic and fine graphs.

\begin{prop}[{\cite{Dthesis}, \cite[\S2.3]{DYaman}}] \label{prop;cones_are_nicer}
In any fine graph, cones are finite.

In any $\delta$-hyperbolic graph, geodesic triangles are conically thin: for any geodesic triangle $([a,b], [b,c], [a,c])$, any edge $e$ on $[a,b]$ is contained in a cone of parameter $50\delta$ around an edge $e'$ that is either in $[a,c]$ at same distance from $a$ than $e$, or in $[b,c]$ at same distance from $b$ than $e$.
\end{prop}

The statement above differs slightly from that in~\cite[\S2.3] {DYaman}, however the proof is the same. Let us present it, as it is very simple. Consider a subsegment $\sigma_1$ of $[a,b]$ that contains $e$, and extends to a distance $3\delta$ from both its end points (unless it reaches $a$ or $c$ before, a case we leave to the reader). Then append, at both ends, at most $2\delta$ long segments $\sigma_0$ and $\sigma_2$ in order to reach $[a,c] \cup [b,c]$ (by usual thin triangles property). If both $\sigma_0$ and $\sigma_2$ reach $[a,c]$, or if both reach $ [b,c]$, then take the arc $\sigma_3$ of that segment, that closes a loop $\sigma_0\sigma_1\sigma_2\sigma_3$. Notice that $\sigma_3$ is of length at most $11\delta$ (and set $\sigma_4, \sigma_5$ to be trivial in order to continue syntactically the argument). If the reached segments are different, close the loop as an hexagon around the center of the triangle, by an arc $\sigma_3$ of length $10\delta$ on $[b,c]$ of the triangle, an arc $\sigma_4$ of at most $2\delta$ between $[a,c]$ and $ [b,c]$, and an arc $\sigma_5$ back to the endpoint of $\sigma_2$ on $[c,a]$, hence of length at most $21\delta$. We thus get a loop $(\sigma_0\sigma_1\sigma_2\sigma_3\sigma_4\sigma_5)$ containing $e$. Triangle inequality forbids $e$ to be in the transitions arcs $\sigma_0, \sigma_2, \sigma_4$. If it is either in $\sigma_1$ or $\sigma_3$ there is actually nothing to prove. Thus, we can extract the simple loop containing $e$. Again triangular inequality ensures some distance between $\sigma_0, \sigma_2, \sigma_4$, and this simple loop must contain an edge $e'$ of $\sigma_2$ or $\sigma_4$ at same distance from respectively $a$ or $b$ than $e$. The length of the loop being at most $42\delta$, its maximal angle is bounded by this quantity. This makes $e$ belong to the cone of parameter $42\delta$ around $e'$, and the proposition is proved.

The above proposition allows us to deduce that intervals are finite in coned-off graphs, even if those graphs are not locally finite.

\begin{prop}\label{prop;geod_finite}
In any hyperbolic fine graph, between any pair of points $a,x$ the set of points in a geodesic between $a$ and $x$ is finite.
\end{prop}

A more important use of cones was explored in~\cite{Israel}, and follows from the fact that in a hyperbolic graph, quasi-geodesics ``without detours'' are conically close to geodesic (\cite[Prop.~1.11]{Israel}). We will use a variant of it, exposed in the next section.

\subsection{Conical \texorpdfstring{$\alpha$}{alpha}-geodesics}

In order to define a flow in the next section, we need to be able to travel coarsely from a point toward another point in a uniformly finite way. In~\cite{AL}, the authors use the concept of $\alpha$-geodesic, for some $\alpha\geq 0$. Precisely, in a metric space, given two points $x$ and $a$, we say that $t$ belongs to an $\alpha$-geodesic between $x$ and $a$ if
\[
d(x,t)+d(t,a)\leq d(x,a)+\alpha.
\]
Note that being on a $0$-geodesic is being on a geodesic. However, when the space is not uniformly locally finite, the set of points in a $\alpha$-geodesic between $x$ and $a$ can be infinite in general, if $\alpha \geq 1$, even if we assume the graph to be fine. For instance, on the coned-off graph of a relatively hyperbolic group, when a geodesic passes through a cone point, any point on the coset of that cone point will be on a $2$-geodesic between the endpoints of that geodesic. But the coset points that are close (in the coset metric) to the entrance or exit points of the geodesic in the coset, will make small angles with some edges of the geodesic, whereas the other ones will make a large angle with any edge of the geodesic. We will use a variant of the notion of $\alpha$-geodesics, that takes angles into account. Constants are not really relevant, but they need to be fixed by convention.

\goodbreak
\begin{defi}\label{defi;slice}
Let $X$ be a graph, and $x,a\in X^{(0)}$. For $\rho\geq 0$ we denote by $\mathcal{E}_{a,x} (\rho)$ the set of edges $e$ with $d(a,o(e)) = \rho$ that are contained in geodesics from $a$ to $x$, or from $x$ to $a$. We say that $t\in X^{(0)}$ \emph{belongs to an $\alpha$-conical-geodesic} between $x$ and $a$ if $t$ is on an $\alpha$-geodesic, $d(a,t)\leq d(a,x)$, and
\[
t\,\in\, \bigcap_{e\in \mathcal{E}_{a,x}(d(a,t)) }
\Co_{40\alpha} (e).
\]
In other words, points in an $\alpha$-conical-geodesic are off a geodesic only by a small angle. We will denote $U_{\alpha}[a,x]$ the set of points belonging to an $\alpha$-conical-geodesic between $x$ and $a$.

For each $\rho\geq 0$, we denote by $\mathcal{S}_{a,x}(\rho)$ the \emph{slice of $U_{\alpha}[a,x]$ at distance $\rho$ from $a$}, that is the set
\[
\mathcal{S}_{a,x}(\rho)=\{t\in U_{\alpha}[a,x]\,|\,d(t,a)=\rho\}
\]
of points in $U_{\alpha}[a,x]$ at distance $\rho$ from $a$.
\end{defi}

We will now see some useful properties, namely that these sets are finite, stable (analogue of~\cite[3.3]{AL} and~\cite[3.4]{AL}), non-empty and slices at large angles are reduced to a point.

\begin{prop}\label{prop;recap}
Let $X$ be a $\delta$-hyperbolic fine graph, and let $a,x \in X$
\begin{enumerate}
\item (Finiteness)\label{lem;slices_are_finite}
For any $\alpha>0$ the set $U_\alpha[a,x]$ is finite. Each slice is contained in a cone of parameter $40\alpha$.
\item (Stability)\label{lem;stab}
If $\alpha =2\delta$, and if $u\in U_{4\delta}[a,x]$ and $v\in U_{2\delta}[a,u]$ and if $d(u,v)\geq 5\delta$ then $v\in U_{4\delta}[a,x]$.
\item\label{lem;slices_are_nonempty}
For all $v$ on a geodesic between $a$ and $x$, the slice of $ U_{2\delta}[a,x]$ at distance $d(a,v)$ from $a$ contains $v$, hence is non empty.
\item\label{lem;un_goulet_de_slice}
If $\A_c(a,x) >(1000\delta)^2$, and $\alpha =2\delta$, then $\mathcal{S}_{a,x} (d(a,c))= \{c\}$.
\end{enumerate}
\end{prop}

\begin{proof}
\begin{step}[\eqref{lem;slices_are_finite}]
By definition, $U_{\alpha}[a,x]$ is contained in an intersection of cones, which are finite according to Proposition~\ref{prop;cones_are_nicer}. Each slice, by definition, is contained in $U_{\alpha}[a,x]$, which is an intersection of a family of cones of parameter $40\alpha$.
\end{step}

\begin{step}[\eqref{lem;stab}]
That $v$ satisfies $d(a,v) + d(v,x) \leq d(a,x)+4\delta$ is~\cite[3.4]{AL}. Consider $y$ in a geodesic $[a,u]$ such that $d(a,y) = d(a,v)$, and an edge $e$ of $[u,a]$ with origin $y$. Then $v\in \Co_{80\delta}(e)$.

We now apply conical thinness of triangles (Proposition~\ref{prop;cones_are_nicer}), to the triangle $([a,x], [x,u], [u,a])$, for the edge $e\in [a,u]$. Since $e$ is at least $5\delta$-far from $u$, the proposition ensures that there exists an edge $e'$ of $[a,x]$ at same distance from $a$ than $e$, for which $e\in \Co_{50\delta} (e')$. It follows, by Proposition~\ref{prop;composition_of_cones}, that $v\in \Co_{130\delta} (e')$, which establishes the claim.
\end{step}

\begin{step}[\eqref{lem;slices_are_nonempty}]
Call $[a,x]_1$ the given geodesic containing $v$. The claimed statement amounts to check that, if $[a,x]_2$ is any geodesic between $a$ and $x$, and $e$ is an edge of $ [a,x]_2$ starting at distance $d(a,v)$ from $a$, then $v$ is in $\Co_{80\delta}(e)$. This is an application of the conical thinness of geodesic triangles, Proposition~\ref{prop;cones_are_nicer} to the degenerate triangle $([a,x]_1, [x,x], [x,a]_2)$ (where $[x,a]_2$ is $[a,x]_2$ with reverse orientation).
\end{step}

\begin{laststep}[\eqref{lem;un_goulet_de_slice}]
Consider two edges $e,e'$ at $c$ (i.e. $o(e)=o(e')=c$) on a geodesic $[a,x]$, where $t(e) $ is closer to $a$ and $t(e')$ is closer to $x$. By definition of slices, one has $ \mathcal{S}_{a,x} (d(a,c)) \subseteq \Co_{80\delta}(e) \cap \Co_{80\delta}(e') \cap \{t, d(c,t)\leq 20\delta \}$. Note that it implies $ \mathcal{S}_{a,x} (d(a,c)) \subseteq \Co_{80\delta+1}(\bar{e}) \cap \Co_{80\delta+1}(\bar{e'}) \cap \{t, d(c,t)\leq 20\delta\}$. We will prove that this intersection is reduced to $\{c\}$.

By Proposition~\ref{prop;goulet}, $\A_c(e, e') >(999\delta)^2$. We apply Lemma~\ref{lem;thetasquare} for $\theta=200\delta$. For all $v\in \Co_{200\delta}(\bar e)$ different of $c$, but at distance at most $20\delta$ from it, one has $\A_c(o(\bar e), v) \leq ((200\delta)^2 + 600\delta)/2$. If $v$ is in $\Co_{200\delta}(\bar{e'})$ as well, then $\A_c(o(\bar e), v) \leq ((200\delta)^2 + 600\delta)/2$. By triangular inequality of angles at $c$, $\A_c(\bar e, e')\leq (200\delta)^2 + 600\delta$, and this contradicts the previous lower bound of $(999\delta)^2$. Thus, the intersection of $\Co_{200\delta}(\bar e)$, $\Co_{200\delta}(\bar {e'})$ and of the ball of radius $20\delta$ around $c$ is reduced to $\{c\}$.

Thus, only $c$ can be in the slice $\mathcal{S}_{a,x} (d(a,c))$. Since the slice is non-empty, we have the result.
\end{laststep}
\let\qed\relax
\end{proof}

\section{The flow on the coned-off graph}

Let us recall that, given a set $Y$ and an additive semigroup $V$, a \emph{flow} is a map $\mathscr{F}: Y\times V \to Y$ such that for all $y$, $\mathscr{F}(y,0)=y$ and $\mathscr{F}(y, t_1+t_2) = \mathscr{F}(\mathscr{F}(y, t_1), t_2)$. In this definition, the semi-group $V$ is usually $\mathbb{R}_+$, but we will use $\mathbb{N}$ (by convention $0\in \mathbb{N}$), and in this case the flow is entirely determined by the \emph{flow step}, which is the map $\mathscr{T}: Y\to Y$ defined by $\mathscr{T} (y) = \mathscr{F}(y,1) $.

In the following, $X$ will be a $\delta$-hyperbolic fine graph (for some integer $\delta >0$), and $\Proba(X)$ is the space of probability measures supported on $X^{(0)}$. We will define a flow for the semi-group $\mathbb{N}$, on the set $Y= \Proba(X) \times X^{(0)}$, which is equivariant by the group of isometries of $X$. We think of it as a version of a geodesic flow.

When $X$ is countable (which is the case in our context) elements of $\Proba(X)$ are just functions from $X^{(0)}$ to $\mathbb{R}_+$ of sum $1$, but we find it convenient and enlightening to see them as probability measures.
\subsection{The flow step}\label{sec;def_flowstep}

In this subsection we will define the flow step
\[
T: \left\{
\begin{aligned} \Proba(X) \times X^{(0)} & \to \Proba(X) \times X^{(0)} \\
(\eta, a) & \mapsto (T_a(\eta), a)
\end{aligned} \right.
\]

First let us introduce some vocabulary. Given $a$ and $x$, we say that the step $T_a$ at $x$ is
\begin{itemize}
\item \emph{initial} if $d(a,x)>5\delta$, and $d(a,x)$ is not a multiple of $5\delta$,
\item \emph{regular} if $d(a,x)>5\delta$ and $d(a,x)$ is a multiple of $5\delta$,
\item \emph{ending} if $a$ is of infinite valence and $1<d(a,x)\leq 5\delta$, or if $a$ has finite valence, $d(a,x) \leq 5\delta$ and there exists $c$ such that $\A_c(a,x) > 1000\delta$,
\item \emph{stationary} in all other cases, i.e. in any of the following three cases
\begin{itemize}
\item[$\bullet$] $x=a$,
\item[$\bullet$] $a$ is of infinite valence and $x$ is a neighbor $a$,
\item[$\bullet$] $a$ is of finite valence, $d(a,x) \leq 5\delta$ and for all $c$, $\A_c(a,x) \leq1000\delta$.
\end{itemize}
\end{itemize}
We denote by
$\delta_x $ the Dirac mass at $x$ and define the flow step $T$ by defining $T_a(\delta_x)$, for $a, x$ two vertices of $X$ and then extend by linearity on probability measures:
\[
T_a\left(\sum_x \lambda_x \delta_x\right) = \sum_x \lambda_x T_a(\delta_x).
\]

For the following definition, we set the parameter $\alpha$ of the definition of slices, to be $\alpha = 2\delta$.

\begin{defi}\label{defi:flowstep}
Let $X$ be a $\delta$-hyperbolic fine graph, and let $a,x\in X^{(0)}$. We set $r_{a,x}$ to be the largest integer $r$
such that $5\delta\,r < d(a,x)$.
\begin{itemize}
\item If the step $T_a$ at $x$ is either initial or regular, then $T_a(\delta_x)$ is the uniform probability measure supported by the slice $\mathcal{S}_{a,x}=\mathcal{S}_{a,x}(5\delta\,r_{a,x})$ (as in Definition~\ref{defi;slice}).
\item If the step $T_a$ at $x$ is ending, and $a$ has infinite valence then $T_a(\delta_x)$ is the uniform probability measure supported by the slice $\mathcal{S}_{a,x} = \mathcal{S}_{a,x}(1)$ at distance $1$ from $a$. If the step is ending and $a$ has finite valence, then $T_a(\delta_x)$ is the Dirac mass supported on the unique $c$ minimizing $d(a,c)$ such that $\A_c(a,x) > 900\delta$.
\item If the step $T_a$ at $x$ is stationary, then $T_a(\delta_x) =
\delta_x$ (and $\mathcal{S}_{a,x} = \{x\}$).
\end{itemize}
\end{defi}

Let us check that this definition makes sense. The set $\mathcal{S}_{a,x}$ is well defined: in the case of an ending step, any vertex $c$ such that $\A_c(a,x) > 900\delta$ is on every geodesic between $a$ and $x$, thus there is a unique one that is closest to $a$. Then the set $\mathcal{S}_{a,x}$ is never empty by Proposition~\ref{prop;recap}$\MK$\eqref{lem;slices_are_nonempty} and it is always finite, by part~\eqref{lem;slices_are_finite} of that same proposition. Observe that among iterates of the flow step applied at $(\delta_x, a)$, only the first iteration can be initial, only one step among the iterates can be ending, and after the ending step, all steps are stationary.

This defines uniquely the value of $T_a(\delta_x)$ for all $x$ and all $a$ and so completely defines the flow step $T$, hence also the flow
\begin{align*}
\mathscr{F}:\Proba (X) \times X^{(0)}\times\mathbb{N}&\to\Proba (X) \times X^{(0)}\\
(\eta,a,k)&\mapsto (T^k_a(\eta),a).
\end{align*}

\subsection{Stationarity and mask}

For two points $x$ and $a$ in $X^{(0)}$, the flow step $T_a(\delta_x)$ is a probability measure whose support is closer to $a$ than $\delta_x$, so the flow step has pushed the point $x$ towards $a$. Iterating this flow step will eventually stall and give a measure that is close to $a$.

\begin{prop}
For $X$ a $\delta$-hyperbolic fine graph and all $a,x\in X^{(0)}$, the flow step from Definition~\ref{defi:flowstep} is stationary in $k$: if
$k>r_{a,x}+1$, one has $T_a^{r_{a,x} +1} (\delta_x)=T_a^{k} (\delta_x)$.
\end{prop}
\begin{proof}
If $y, y'$ are in the support of $T_a(\delta_x)$, one has $d(a,y) = d(a,y') \leq d(a,x)$. In case of equality, then $T_a(x) = \delta_x$. Indeed, if the flow step is initial or regular, then $d(a,y) = d(a,y') < d(a,x)$, since by definition, $5\delta\,r_{a,x} < d(a,x)$. The ending step case, also induces a strict inequality, and the stationary step case induces the equality $T_a(x) = \delta_x$. Hence the flow eventually become stationary for $k$ large enough.
\end{proof}

\begin{defi}\label{flot}
Given any two vertices $a$ and $x$ of $X$ a $\delta$-hyperbolic fine graph, the \emph{mask of $a$ for $x$} is the probability measure given by
\[
\mu_x(a)=T^{R(a,x)}_a (\delta_x) \, =\, \lim_{k\to\infty} T^k_a (\delta_x),
\]
where $R(a,x)$ the minimal number of iterations of $T_a$ on $\delta_x$ so that the stationary step is reached.
\end{defi}

Let us observe the following.

\begin{prop}\label{prop;theta0}
Let $X$ be a $\delta$-hyperbolic and fine graph. For $a,x\in X^{(0)}$, then
\begin{enumerate}
\item \label{2.4.1} If $a$ is of infinite valence, for all $x$, the measure $\mu_x(a)$ is supported by the $1$-neighborhood of $a$. Moreover, if $e$ is the first edge of a geodesic segment $[a,x]$, the measure $\mu_x(a)$ is supported in ${\Co}_{160\delta}(e)$.
\item \label{2.4.2} If $a$ is of finite valence, for all $x$, the measure $\mu_x(a)$ is supported by ${\Co}_{\theta_0}(a)$ for $\theta_0=1160\delta$.
\end{enumerate}
\end{prop}

The last estimate is very crude, and with some closer look, one is likely to get a much more precise control.

\begin{proof}
\begin{step}[\eqref{2.4.1}] If $a$ has infinite valence, by definition of the ending step in this case the measure is in the slice at distance 1 from $a$. By Proposition~\ref{prop;recap}$\MK$\eqref{lem;stab}, the last slice used is in a $4\delta$-conical geodesic from $a$ to $x$. So, given $e$ which is first edge on a geodesic from $a$ to $x$, by the definition of slices, this last slice used is contained in ${\Co}_{160\delta}(e)$.
\end{step}

\begin{laststep}[\eqref{2.4.2}]
In the case of a finite valence vertex $a$, the condition for the step being stationary is that there are no angle larger than $1000\delta$ between $a$ and a point in the support. Again, by Proposition~\ref{prop;recap}$\MK$\eqref{lem;stab}, the support is contained in a slice is in a $4\delta$-conical geodesic from $a$ to $x$, hence in a cone $\Co_{160\delta} (e_0)$ centered at an edge in some geodesic $[a,x]$, at distance at most $5\delta$ from $a$. We moreover know, by definition of the ending step, that there is no angle larger than $1000\delta$ between $a$ and $o(e_0)$. Thus, by composition of cones, Proposition~\ref{prop;composition_of_cones}, the support is contained in $\Co_{1160\delta} (e)$ for an initial edge of a geodesic $[a, x]$.
\end{laststep}
\let\qed\relax
\end{proof}

The definition of the flow step being in a purely metric way, as a consequence the masks are equivariant with respect to the isometry group.

\begin{prop}\label{prop;mask_equivariance}
Let $X$ be a $\delta$-hyperbolic fine graph.
The map
\[
X^{(0)}\times X^{(0)}\to\Proba (X),\ (x,a)\mapsto \mu_{x}(a)
\]
that associates to a pair of points $(x,a)$ the mask of $a$ for $x$ is equivariant with respect to the isometry group of $X$. More precisely, if $\phi$ is an isometry of $X$, one has $\phi_*\mu_{x}(a) = \mu_{\phi (x)} (\phi(a))$.
\end{prop}

\subsection{The flow stays focused}

In this subsection we will be recording the behavior of the flow step at different stages

\begin{lemm}\label{focusLemma}
Let $X$ be a $\delta$-hyperbolic fine graph, and for $a\in X^{(0)}$ let $T_a$ denote the flow step as in Definition~\ref{defi:flowstep}. Let $x,x'\in X^{(0)}$.
\begin{enumerate}
\item \label{lem;stayfocus_regular}
If $x$ is such that the $k$ first iterates of $T_a$ are initial or regular steps, then, for any two points $y,y'$ in the support of $T^k_a(\delta_x)$, one has $d(y,y') \leq 5\delta$. Moreover the support of $T^k_a(\delta_x)$ is contained in a cone of parameter $160\delta$.
\item \label{lem;focus_initial}
If $x,x'$ are neighbors, and $5k\delta<d(a,x)\leq d(a,x') < 5(k+1)\delta$, then the union of the supports of $T_a(\delta_x)$ and $T_a(\delta_{x'})$ has diameter at most $8\delta+1$.
\item \label{lem;elementary_confluence}
If $d(x,x')\leq 8\delta$, with $d(x,a) = d(x',a)$ and assume that for both of them, the flow step is regular. Then the union of the supports of both measures $T_a(\delta_x)$ and $T_a(\delta_{x'})$ has diameter at most $8\delta$. Moreover, the union of these supports is contained in a cone of parameter $170 \delta$, and their intersection is not empty.

\item \label{prop;firststep}
If $x,x'$ are neighbors and both at distance larger than $5\delta$ from $a$, then there are two exponents $\epsilon,
\epsilon'$ that are either $1$ or $0$, such that $T^\epsilon_a(\delta_x)$ and $T^{\epsilon'}_a(\delta_{x'})$ have support on a sphere centered at $a$ of radius $5k\delta$ for some $k$, and the diameter of their union is at most $8\delta+2$.
\end{enumerate}
\end{lemm}

\begin{proof}
\begin{step}[\eqref{lem;stayfocus_regular}]
By Proposition~\ref{prop;recap}$\MK$\eqref{lem;stab} $y$ and $y'$ are on $U_{4\delta}[a,x]$. As in~\cite{AL} we can estimate:
\[
d(a,x) + d(y,y') \leq \max \{ d(a,y) + d(y',x),
\, d(a,y') + d(y,x) \} +\delta \; \leq d(a,x) + 4\delta +\delta.
\]
Hence $d(y,y') \leq 5\delta$. Then, observe that by Proposition~\ref{prop;recap}$\MK$\eqref{lem;stab} the support of $T^k_a(\delta_x)$ is in a slice of $U_{4\delta} [a,x]$. Proposition~\ref{prop;recap}$\MK$\eqref{lem;slices_are_finite} concludes.
\end{step}

\begin{step}[\eqref{lem;focus_initial}]
By assumption the steps of the flow are initial. Let $y_0$ and $y'_0$ be respectively on geodesics $[x,a]$ and $[x',a]$ at distance $5k\delta$ from $a$. By $\delta$-hyperbolicity, they are at distance less than $2\delta+1$ from each other. Let $y, y'$ respectively be in the supports of $T_a(\delta_x)$ and $T_a(\delta_{x'})$. Apply the argument of part~\eqref{lem;stayfocus_regular} using that $y$ and $y'$ are on $2\delta$-geodesics (there is only one step of flow). One obtains that $d(y,y_0) \leq 3\delta$ and $d(y', y_0') \leq 3\delta$, hence $d(y,y')\leq 8\delta+1$.
\end{step}

\begin{step}[\eqref{lem;elementary_confluence}]
Here $x$ and $x'$ are not assumed neighbors, but at distance $5(k+1)\delta$ from $a$. Let $y_0$ and $y'_0$ respectively be on geodesics $[x,a]$ and $[x',a]$ at distance $5k\delta$ from $a$. By $\delta$-hyperbolicity, they are at distance less than $2\delta$ from each other and the same argument than part~\eqref{lem;focus_initial} gives a bound of $8\delta$ on the union of the supports of $T_a(\delta_x)$ and $T_a(\delta_{x'})$

Given two geodesics $[a,x]$ and $[a,x']$, we denote by $e, e'$ the edges of these geodesics starting at distance $\left(d(a,x) -5\delta\right)$ from $a$. By $\delta$-hyperbolicity, the two segments stay $2\delta$-close for a length of at least $\delta$ after $e$ and $e'$. One deduces (as in the proof of Proposition~\ref{prop;goulet}) that $e\in \Co_{10\delta}(e')$. Therefore, by composition of cones (Proposition~\ref{prop;composition_of_cones}) $\Co_{160\delta}(e) \subseteq \Co_{170\delta} (e')$. Moreover, $o(e)$ is at distance at most $\delta$ from $[a,x']$, so is on a $2\delta$-geodesic between $a$ and $x'$. Also, as we discussed, it is in $\Co_{10\delta}(e')$, for any edge $e'$ on any geodesic $[a,x']$ satisfying that $e'$ starts at $\left(d(a,x) -5\delta\right)$ from $a$. Thus, $o(e) \in U_{2\delta}[a,x']$ and hence by Proposition~\ref{prop;recap}$\MK$\eqref{lem;slices_are_nonempty} the intersection of the supports is non-empty.
\end{step}

\begin{laststep}[\eqref{prop;firststep}]
There are several cases, depending whether the flow from $x$, respectively from $x'$, needs an initial step or not, in other words, whether $x$, respectively $x'$, are at distance strictly greater than $5k\delta$ or equal to $5k\delta$ from $a$.

If the first step of the flow is initial for both $(\delta_x, a)$ and $(\delta_{x'}, a)$, then part~\eqref{lem;focus_initial} says that the union of the supports of $T_a(\delta_x)$ and $T_a(\delta_x)$ has diameter at most $8\delta$. The same is true according to part~\eqref{lem;elementary_confluence} if both first steps of the flow are regular.

Assume that $d(a,x) = 5k\delta$ and $d(a,x') >5k\delta$. Then we compare $\delta_x$ (i.e. $\epsilon =0$) with $T_a(\delta_{x'})$ (i.e. $\epsilon' =1$). According to part~\eqref{lem;stayfocus_regular} the support of $T_a(\delta_{x'})$ has diameter at most $5\delta$, and by Proposition~\ref{prop;recap}$\MK$\eqref{lem;slices_are_nonempty} it contains a point at distance $1$ from $x'$, hence at distance less than $2$ from $x$, the union of the supports of $T^\epsilon_a(\delta_x)$ and $T^{\epsilon'}_a(\delta_{x'})$ has diameter at most $5\delta+2$, hence the result.
\end{laststep}
\let\qed\relax
\end{proof}

\subsection{Check points at large angles}

The following proposition says that when the flow passes a large angle, this large angle acts as a check point: The flow could have started at this large angle, and still give the same result.

\begin{prop}\label{prop;checkpoint}
Let $(X,d)$ a $\delta$-hyperbolic fine graph. Let $a, x$ be vertices in $X$. If there exists $c\in [a,x]$ such that $\A_c(a,x)>(2000 \delta)^2$, then for any $x'$ at distance~$1$ from $x$, we have that
\[
\mu_x(a) = \mu_{x'}(a)= \mu_c(a).
\]
\end{prop}

\begin{proof}
Define $r(c)$ to be the largest $r$ so that $(r_{a,x} -r)5\delta\geq d(a,c)$. First notice that for all $r\leq r(c)$, and all $b$ in the support of $ T^r_a(\delta_x) $, one has
\[
\A_c(a,b)>(1910\delta)^2.
\]
Indeed, let $b$ in the support of $ T^r_a(\delta_x) $. Then $d(a,b) = (r_{a,x} -r)5\delta >d(a,c)$, and by Proposition~\ref{prop;recap}$\MK$\eqref{lem;slices_are_nonempty} there is $b'$ on the geodesic $[a,x]$ realizing the maximal angle at $c$, that is in the support of $ T^r_a(x) $. Hence $\A_c(a,b') >(2000\delta)^2$. Since $b'$ and $b$ are in a same cone of parameter $160\delta$ (Proposition~\ref{prop;recap}$\MK$\eqref{lem;slices_are_finite}), the maximal angle between them is at most $2(160\delta)^2$ (by Lemma~\ref{lem;thetasquare}). Thus, $\A_c(a,b)>(1987\delta)^2>(1910\delta)^2$ as claimed.

We now show that $T^{r(c)+1}_a(\delta_x)=T_a(\delta_c)$. If $d(a,c)$ is a multiple of $5\delta$, for any point $b$ in the support of $ T^{r(c)-1}_a(\delta_x)$ then $\A_c(a,b)>(1910\delta)^2$. Thus, all geodesics from $a$ to $b$ pass through $c$ and make an angle at least $(1900\delta)^2$. It follows that the intersection of all cones of parameter $80\delta$ centered at an edge of $\mathcal{E}_{a,b}(d(a,c))$ is reduced to $\{c\}$ (by Proposition~\ref{prop;recap}$\MK$\eqref{lem;un_goulet_de_slice}), hence $T_a(\delta_b) = \delta_c$, which means that $T^{r(c)}_a(\delta_x) = \delta_c$, and $T^{r(c)+1}_a(\delta_x)=T_a(\delta_c)$ as claimed. When $d(a,c)$ is not a multiple of $5\delta$, apply $T_a$ to all points $b$ in the support of $ T^{r(c)}_a(\delta_x) $. But since we saw that $\A_c(a,b)>(1910\delta)^2$, all geodesics from $a$ to $b$ pass through $c$ and make an angle at least $(1900\delta)^2$. The points on $2\delta$-geodesics from $a$ to $c$ are exactly the points on $2\delta$-geodesics from $a$ to $b$ that are at distance at most $d(a,c)$ from $a$, and we have equality of the sets $\mathcal{E}_{a,b}((r_{a,x}- r(c))5\delta)=\mathcal{E}_{a,c}((r_{a,x}- r(c))5\delta)$. Hence $T_a(\delta_b) = T_a(\delta_c)$, and $T^{r(c)+1}_a(\delta_x)=T_a(\delta_c)$ again as claimed.

Finally, if $x'\neq c$ and $x'$ at distance $1$ from $x$, then, for $r'(c)$ defined similarly to $r(c)$ for $x'$ we have
\[
T^{r(c)+1}_a(\delta_x) = T^{r'(c)+1}_a(\delta_{x'}) = T_a(\delta_c).
\]
The angle at $c$ between $a$ and $x'$ can be reduced by one unit, which is still well above the threshold to apply the previous argument. Iterating the flow step gives the conclusion of the proposition.
\end{proof}

\subsection{Confluence}

We will need to control that, given $a$, and $x,x'$ close to each other, but ``far'' from $a$ (in terms of distance or of angles), each of the regular flow steps from $(\delta_x, a)$ and from $(\delta_{x'}, a)$, are uniformly close. To do that, we will be talking of confluence. Recall that the norm $\|\eta\|$ of a measure is defined as its $\ell^1$-norm $\|\eta\| = \sum_{x\in X} |\eta(x)|$. For $\eta, \eta'$ two probability measures on $X$, denote by $\mathscr{M}(\eta,\eta')$ the measure defined by
\[
\mathscr{M}(\eta,\eta') (\{y\}) = \min \{ \eta (\{y\}), \eta'(\{y\}) \}.
\]
We also define the \emph{symmetric difference} of $\eta$ and $\eta'$ as
\[
(\eta \Delta\eta') = \eta+\eta' - 2\mathscr{M}(\eta,\eta'),
\]
which is always a positive measure since $\eta, \eta'$ are probability measures. Moreover, $\|\eta \Delta\eta'\| = \| \eta - \eta'\|$.

\begin{defi}\label{defconf}
Let $\beta \geq 0$. We say that the flow step $T_a$ is \emph{$\beta$-confluent} on $(\eta,
\eta')$ if
\[
\| \mathscr{M} (T_a (\eta), T_a(\eta')) \| \geq
\| \mathscr{M} (\eta, \eta') \| + \beta \| \eta \Delta
\eta'\|.
\]
\end{defi}
\begin{rema}\label{betterconf}
Notice that the flow step $T_a$ is $\beta$-confluent on $(\eta,\eta')$ if and only if
\[
\| T_a(\eta) \Delta T_a(\eta') \| \leq (1- 2\beta) \| \eta \Delta
\eta' \|.
\]
Indeed, observe that $\| T_a(\eta) \Delta T_a(\eta') \| = 2- 2 \| \mathscr{M} (T_a (\eta), T_a(\eta'))
\|$, and similarly $\| \eta \Delta
\eta' \| = 2 - 2 \mathscr{M} (\eta, \eta')$. After substitution, and dividing by $2$, the above condition reads:
\[
1 - \| \mathscr{M} (T_a (\eta), T_a(\eta'))
\| \leq (1 - \|\mathscr{M} (\eta, \eta') \|) - \beta \| \eta \Delta
\eta' \|
\]
which is equivalent to the inequality of the
definition.
\end{rema}

\begin{prop}\label{prop;decay_in_C}
Let $X$ be a fine $\delta$-hyperbolic graph with a co-finite group action. Let $k\geq 1$, and $\eta, \eta'$ be two measures whose supports are on $S(a,5(k+1)\delta)$, and have union of diameter less than $8\delta$.

Then, there is a constant $C>1$ such that the $k$ consecutive flow steps on $(\eta, a)$ and $(\eta',a)$ are $\frac{1}{C}$-confluent. In particular, for all $m\geq k$,
\[
\| T^m_a(\eta) \Delta
T^m_a(\eta') \| \leq \left(1-\frac{2}{C}\right)^k \| \eta \Delta
\eta' \|.
\]
\end{prop}

\begin{proof}
First notice that the assumption that $X$ is a fine $\delta$-hyperbolic graph with a co-finite group action implies that here are only finitely many orbits of cones of given parameter, and in each of these orbits, cones are finite (Proposition~\ref{prop;cones_are_nicer}), with same cardinality. Hence, there exists $C>1$ that is an upper bound on the cardinality of each cone of parameter $160\delta$ in $X$.

We first proceed by induction to prove that for each $i\leq k$, the union of supports of the measures $T^i_a(\eta)$ and $T^i_a(\eta')$ have diameter at most $8\delta$. There is nothing to prove for $i=0$. For $i>0$, assuming $T^{i-1}_a(\eta)$ and $T^{i-1}_a(\eta')$ have support whose union has diameter at most $8\delta$, we can say that, since all steps of the flow so far were regular, these supports are on the sphere $S(a, 5 (k+2-i)\delta)$. We apply Lemma~\ref{focusLemma}$\MK$\eqref{lem;elementary_confluence} to obtain that the union of supports of $T^{i}_a(\eta)$ and $T^{i}_a(\eta')$ have diameter at most $8\delta$.

We then prove the confluence. For each Dirac masses $\lambda_v\delta_v$ and $\lambda_{v'}\delta_{v'}$, respectively in the support of $T^{i-1}_a(\eta) - \mathscr{M} (T^{i-1}_a (\eta), T^{i-1}_a(\eta')) $ and of $T^{i-1}_a(\eta')- \mathscr{M} (T^{i-1}_a (\eta), T^{i-1}_a(\eta')) $, the regular step of the flow sends each of these Dirac masses on a uniform measure on a subset of a slice, in a cone of parameter $160\delta$, by Proposition~\ref{prop;recap}$\MK$\eqref{lem;slices_are_finite}. Moreover the respective supports of these uniform measures share at least one point, by Lemma~\ref{focusLemma}$\MK$\eqref{lem;elementary_confluence}. Therefore, this step is $\frac{1}{C}$-confluent on the Dirac masses $\lambda_v\delta_v$ and $\lambda_{v'}\delta_{v'}$ (see Definition~\ref{defconf} above), where $C$ is the upper bound on the cardinality of cones of parameter $160\delta$ in $X$. Precisely, according to Remark~\ref{betterconf} we have that
\[
\| T_a(\delta_v) \Delta T_a(\delta_{v'}) \| \leq \left(1- \frac{2}{C}\right) \|\delta_v \Delta\delta_{v'} \|.
\]
Since the flow step on measures is defined by linearity on Dirac masses, this confluence on all masses in these supports gives that
\[
\|T^{i}_a(\eta) \Delta T^{i}_a(\eta')\| = \|T(T^{i-1}_a(\eta)) \Delta T(T^{i-1}_a(\eta'))\| \leq \left(1-\frac{2}{C}\right) \| T^{i-1}_a(\eta)
\Delta T^{i-1}_a(\eta') \|\qedhere
\]
\end{proof}

\begin{prop}\label{prop;kappa}
Let $X$ be a hyperbolic fine graph, with a cofinite group
action. There exists a constant $\kappa <1$ such that, for all $a$, for all $x_1, x_2$ at distance $1$ from each other, if $x$ is among $x_1,x_2$, closest to $a$, and if $\Theta$ denotes the minimal sum of angles at infinite valence vertices on a geodesic $[a,x]$, then for $d(a,x)+\Theta $ sufficiently large, one has
\[
\| \mu_{x_1}(a) \Delta \mu_{x_2}(a) \| \leq \kappa^{d(a,x)+\Theta}
\]
\end{prop}

Observe that the statement would hold also for $\Theta$ defined as the sum of all consecutive angles, and also if it was taken to be the maximum of these sums.

\begin{proof}
If one of the angles on $[a,x]$ is larger than $(2000\delta)^2$ then by Proposition~\ref{prop;checkpoint}, $\mu_{x_1}(a) \Delta \mu_{x_2}(a) =0$. If all the angles on $[a,x]$ are smaller that $(2000\delta)^2$,
then
\[
d(a,x) \leq d(a,x)+\Theta \leq (2000\delta)^2 d(a,x).
\]

It follows that $(\kappa)^{d(a,x)+\Theta} \geq (\kappa^{(2000\delta)^2 6\delta })^{d(a,x)/6\delta}$. We then choose $\kappa<1$ sufficiently close to $1$ so that $\kappa^{(2000\delta)^2 6\delta } > (1-\frac{2}{C})$ from Proposition~\ref{prop;decay_in_C}. If $d(a,x)$ is sufficiently large (uniformly), $d(a,x)/6\delta$ is larger than the number of regular step of the flow applied to $(\delta_x, a)$. Thus, Proposition~\ref{prop;decay_in_C}, and Lemma~\ref{focusLemma}$\MK$\eqref{prop;firststep} give the conclusion.
\end{proof}

\subsection{Non-confluence}

In this subsection we will be showing that for two given sources of the flow, any point on a geodesic between those two sources has disjoint masks for each of those sources.

\begin{prop}\label{non-conf}
Let $X$ be a $\delta$-hyperbolic and fine graph, and let $\mu$ be the mask of Definition~\ref{flot}.
\begin{enumerate}
\item \label{2.12.1} If $d(x,x')\geq 10\delta$ and if $a$ is of finite valence, at distance at least $5\delta$ from both $x$ and $x'$, and $a\in [x,x']$ for some geodesic, then $\| \mu_x(a) \Delta \mu_{x'}(a) \| =2$.
\item \label{2.12.2} If $a$ has infinite valence, and $\A_a(x,x') > 10(160\delta)^2 $ then $\| \mu_x(a) \Delta \mu_{x'}(a) \| =2$.
\end{enumerate}
\end{prop}

\begin{proof}
\begin{step}[\eqref{2.12.1}]
We proceed by contradiction. Assume that $v$ is in the support of both $\mu_x(a)$ and $\mu_{x'}(a) $.

In the first case, assume that no step in the flow was \emph{ending}. Then $v$ is at distance $5\delta$ from $a$ and is in $4\delta$-geodesics between both $a,x$ and $a,x'$ (by Proposition~\ref{prop;recap}$\MK$\eqref{lem;stab}). This is not possible.

Assume that one step of the flow for $\delta_x$ was ending. If the there was no step of the flow of type ending for $\delta_{x'}$, then the supports of $\mu_x(a)$ and $\mu_{x'}(a)$ are on different spheres around $a$, so they are disjoint. If there was an ending step for $\delta_{x'}$, then, since $a$ is of finite valence, $\mu_x(a)$ and $\mu_{x'}(a)$ are Dirac masses carried by vertices $y,y'$ on which, respectively, $\A_y(a,x)$ and $\A_y(a,x')$ are larger than $20\delta$. Since $a$ is on a geodesic from $x$ to $x'$, by Proposition~\ref{prop;goulet} one has $y\neq y'$.
\end{step}

\begin{laststep}[\eqref{2.12.2}]
The assumption indicates that the flow from $\delta_x$ to $a$ and from $\delta_{x'}$ to $a$ has an ending step. It follows from Proposition~\ref{prop;theta0} that the measures $\mu_x(a)$ and $ \mu_{x'}(a)$ have supports contained in two cones of parameter $160\delta$ centered respectively on an edge starting a geodesic $[a,x]$ and starting a geodesic $[a,x']$. Let $e,e'$ be these edges. One has $\A_a(e,e') > 9 (160\delta)^2$. It follows from Lemma~\ref{lem;thetasquare} that the intersection of the cones is reduced to $\{a\}$. Thus the intersection of the supports is empty.
\end{laststep}
\let\qed\relax
\end{proof}

\begin{coro}\label{cor;ALconique}
Let $X$ be the coned-off graph of a finitely generated relatively hyperbolic group $G$, over its family of peripheral subgroups. There exists $p>1$, and $\varepsilon >0$ such that for all $x$
\[
\sum_{a\in X} \| \mu_1(a) \Delta \mu_{x}(a) \|^p
\]
is convergent, and if $x$ is at distance at least $10\delta$ from $1$, the sum is larger than $\varepsilon d(1,x)$.
\end{coro}

\begin{proof}
We identify $G$ with the image of its Cayley graph in $X$. For the first statement, let $\Theta(x,x')$ denote the minimal sum of angles at infinite valence vertices on a geodesic $[x,x']$. Define $d': X\times X \to \mathbb{N}$ as $d'(g_1,g_2) = d_X(g_1,g_2) + \Theta(x,x')$. By Proposition~\ref{prop;pseudo_word_metric}, $d'$ is coarsely larger than a word metric $d_w$ on $G$. In particular, there is $A\geq 1$ and $ \gamma_G >1 $ such that $S_{d'}(R,G) = \{ g\in G, d'(1,g) = R\}$ has cardinality less than $A\gamma_G^R$. Denote by $S_{d'}(R)$ the union of $S_{d'}(R,G) $ and of the vertices $x$ of $X$ of infinite valence such that $d'(1,x)= R$. One can check that its cardinality is at most $A' \gamma_G^R$ for a certain constant $A'$ (related to the maximal order of intersection of two peripheral subgroups).

Recall that we defined a constant $\kappa$ in Proposition~\ref{prop;kappa}. We choose $p$ such that $\kappa^p< 1/\gamma_G$. Then, fixing $x$, we have the bound
\[
\sum_{a\in S_{d'}(R)} \| \mu_1(a) \Delta \mu_{x}(a) \| \leq
\kappa^{p R} A' \gamma_G^{ R}
\]
This defines a summable family of numbers therefore $\sum_{a\in X} \| \mu_1(a) \Delta \mu_{x}(a) \|^p $ is convergent.


The lower bound is given by Proposition~\ref{non-conf} and the fact that there are at least $d(x,x')/2 - 1$ vertices of finite valence between $x$ and $x'$.
\end{proof}


\section{An action that is proper in the coned-off distance}\label{proper_in_CO}

In this section~$X$ is the coned-off Cayley graph of a relatively hyperbolic group $G$, and $\mathscr{V}=\mathbb{R}X^{(0)}$ is the vector space of all functions from $X^{(0)}$ to $\mathbb{R}$ with finite support, that we endow with the $\ell^1$-norm. For all $p >1$ or $p=\infty$, we consider $\mathscr{W}_p= \ell^p(X^{(0)},
\mathscr{V})$. In other words, an element $\omega$ of $\mathscr{W}_p$ is the data, for all $a\in X^{(0)}$, of a vector $\omega_a\in \mathscr{V}$, such that its norm
\[
\| \omega \|_{\mathscr{W}_p} = \left(\sum_a \| \omega_a \|^p_{\mathscr{V}}
\right)^{\frac1p} = \left(\sum_a \left(\sum_v |\omega_a(v)|
\right)^p \right)^{\frac1p}
\]
is finite. Let $ {\mathscr{W}} $ be one of the ${\mathscr{W}}_p$. We denote by $\overrightarrow{{\Isom}} {\mathscr{W}}$ the unitary group of ${\mathscr{W}}$, i.e. the subgroup of linear automorphisms of ${\mathscr{W}} $ that are isometries for the norm. The group $G$ admits an isometric linear action on ${\mathscr{W}} $ by precomposition by isometries of $X$. Let us denote $\pi: G\to \overrightarrow{{\Isom}} {\mathscr{W}} $ this action. We want to promote it into an action by affine isometries that has no fixed point. For that we need a cocycle, as we recall in the next subsection.
\subsection{Generalities on cocycles and actions} Recall that we may see $ {\mathscr{W}} $ as an affine space, and also as an abelian group (for the addition of vectors). The group of affine isometries of (the affine space) ${\mathscr{W}} $ is the semidirect
product
\[
{\Aff \Isom} {\mathscr{W}} = {\mathscr{W}} \rtimes \overrightarrow{{\Isom}} {\mathscr{W}}
\]
for
the natural action of $\overrightarrow{{\Isom}} {\mathscr{W}} $ on ${\mathscr{W}} $. Given a group $G$, any homomorphism $\phi:G\to {\Aff \Isom}{\mathscr{W}}$ has an image in $\vec{\phi}:G\to \overrightarrow{{\Isom}} {\mathscr{W}} $ through the quotient map, and, given this homomorphism $\vec{\phi}$, the homomorphism $\phi$ is characterised by a map $c:G\to {\mathscr{W}} $ recording $c(g) =
\phi(h)(\vec{0}_ {\mathscr{W}})$. Applying the law of the semidirect product reveals that $c$ satisfy the cocycle relation
\[
c(g_1g_2) = c(g_1) + \vec{\phi}(g_1) (c(g_2)),
\]
for all $g_1,g_2\in G$. The cocycle is actually the difference between $\phi$ and a given section of $\vec{\phi}$ in the semi-direct product. Conversely, whenever one has $\vec{\phi}:G\to \overrightarrow{{\Isom}} {\mathscr{W}}$ and $c:G\to {\mathscr{W}}$ satisfying the cocycle relation for $\vec{\phi}$, one can define the map
\[
\phi:G\to {\Aff \Isom} {\mathscr{W}}, g\mapsto \phi(g) = (c(g), \vec{\phi}(g)),
\]
which is a homomorphism.


\begin{defi}\label{proper}
For a finitely generated group $G$, endowed
with a (locally finite) word metric $d$, the homomorphism $\phi$ is called \emph{proper} if $c$ is proper, or equivalently if $\|c(g_i)\|$ goes to infinity for any sequence $(g_i)$ of elements in $G$ for which $d(1,g_i)$ goes to infinity in $\mathbb{R}$.

If $d$ is a left invariant metric obtained from a coned-off Cayley graph over the cosets of subgroups, (or, in other words, a relative word metric), then $\phi$ is relatively proper if $\|c(g_i)\|$ goes to infinity for any sequence $(g_i)$ of elements in $G$ for which $d(1,g_i)$ goes to infinity.
\end{defi}

\subsection{A relatively proper affine isometric action}

We return to our initial context. We have $G$ a relatively hyperbolic group, and $\pi: G\to \overrightarrow{{\Isom}} {\mathscr{W}} $ for our relatively hyperbolic group $G$, which we want to promote to a homomorphism in ${\Aff}{\Isom} {\mathscr{W}} $. We thus need a cocycle for $\pi$. We identify $G$ with the vertices of its Cayley graph $\Cay (G)$ in the coned-off Cayley graph $X$. We define the map $c$ as follows
\[
c: \left\{
\begin{aligned} G &\to \mathscr{W}_{\infty}
\\ g & \mapsto \{a \mapsto
\mu_1(a) -
\mu_g(a)
\}.
\end{aligned} \right.
\]
where $\mu_x(a)$ is the probability measure from Definition~\ref{flot}. Recall that $\mu_1(a) - \mu_g(a)$ has same $\ell^1$-norm as the symmetric difference $\mu_1(a) \Delta\mu_{g}(a)$ because the latter is the absolute value of the former. The following proves the first part of our main result, Theorem~A

\begin{theo}\label{thm;actionAL}
Let $G$ be a relatively hyperbolic group and let $X$ be the coned-off Cayley graph of $G$ with respect to its peripheral subgroups. For $p>1$ large enough, $c$ has its values in $\mathscr{W}_p=\ell^p(X^{(0)},\mathscr{V})$, and is a cocycle for the representation $\pi$ on $\mathscr{W}_p$.

Moreover, the homomorphism $(c,\pi): G \to {\Aff}{\Isom} {\mathscr{W}}_p$ is relatively proper for the metric on $G$ induced by the coned-off Cayley graph. Precisely, there is $\epsilon >0$ and a constant $k_0$, for which for all $g\in G$, $\|g(\vec{0})\|^p
\geq \epsilon d_X(1,g) -k_0$.
\end{theo}

\begin{proof}
That $c$ is a cocycle for $\pi$ is a standard computation using Proposition~\ref{prop;mask_equivariance}. It takes its values in $\ell^p(X^{(0)},\mathscr{V})$, by Corollary~\ref{cor;ALconique}. The estimate on $\|g(\vec{0})\|$ is a consequence of the last part of Corollary~\ref{cor;ALconique}.
\end{proof}
\begin{thebibliography}{MMMMMM}
\bibitem[AL17]{AL} Alvarez, Aurélien; Lafforgue, Vincent Actions affines isométriques propres des groupes hyperboliques sur des espaces \(\ell\) p , Expo. Math., Volume 35 (2017) no. 1, pp. 103-118 \href{https://ahl.centre-mersenne.org/item/AHL_2020__3__35_0/#AL}{Publisher reference}.
\bibitem[Arn15]{Ar} Arnt, Sylvain Spaces with labelled partitions and isometric affine actions on Banach spaces (2015) (\url{https://arxiv.org/abs/1401.0125}) \href{https://ahl.centre-mersenne.org/item/AHL_2020__3__35_0/#Ar}{Publisher reference}.
\bibitem[BdlHV08]{BdlHV} Bekka, Bachir; de la Harpe, Pierre; Valette, Alain Kazhdan’s Property, New Mathematical Monographs, 11, Cambridge University Press, 2008 \href{https://ahl.centre-mersenne.org/item/AHL_2020__3__35_0/#BdlHV}{Publisher reference}.
\bibitem[BFGM07]{BFGM} Bader, Uri; Furman, Alex; Gelander, Tsachik; Monod, Nicolas Property (T) and rigidity for actions on Banach spaces, Acta Math., Volume 198 (2007) no. 1, pp. 57-105 \href{https://ahl.centre-mersenne.org/item/AHL_2020__3__35_0/#BFGM}{Publisher reference}.
\bibitem[Bou16]{Bou} Bourdon, Marc Cohomologie et actions isométriques propres sur les espaces L p , Geometry, topology, and dynamics in negative curvature (London Mathematical Society Lecture Note Series), Volume 425, Cambridge University Press, 2016, pp. 84-106 \href{https://ahl.centre-mersenne.org/item/AHL_2020__3__35_0/#Bou}{Publisher reference}.
\bibitem[Bow12]{BowRH} Bowditch, Brian H. Relatively hyperbolic groups, Int. J. Algebra Comput., Volume 22 (2012) no. 3, 1250016, 66 pages \href{https://ahl.centre-mersenne.org/item/AHL_2020__3__35_0/#BowRH}{Publisher reference}.
\bibitem[Cla36]{Cla} Clarkson, James A. Uniformly convex spaces, Trans. Am. Math. Soc., Volume 40 (1936) no. 3, pp. 396-414 \href{https://ahl.centre-mersenne.org/item/AHL_2020__3__35_0/#Cla}{Publisher reference}.
\bibitem[Dah03]{Dthesis} Dahmani, François Les groupes relativement hyperboliques et leurs bords, Ph. D. Thesis, Université Strasbourg 1 (France) (2003) \url{http://www.theses.fr/2003STR13033} \href{https://ahl.centre-mersenne.org/item/AHL_2020__3__35_0/#Dthesis}{Publisher reference}.
\bibitem[Dah06]{Israel} Dahmani, François Accidental parabolics and relatively hyperbolic groups, Isr. J. Math., Volume 153 (2006), pp. 93-127 \href{https://ahl.centre-mersenne.org/item/AHL_2020__3__35_0/#Israel}{Publisher reference}.
\bibitem[Day41]{Day} Day, Mahlon M. Some more uniformly convex spaces, Bull. Am. Math. Soc., Volume 47 (1941), pp. 504-507 \href{https://ahl.centre-mersenne.org/item/AHL_2020__3__35_0/#Day}{Publisher reference}.
\bibitem[DS05]{DS} Druţu, Cornelia; Sapir, Mark Tree graded spaces and asymptotic cones of groups, Topology, Volume 44 (2005) no. 5, pp. 959-1058 \href{https://ahl.centre-mersenne.org/item/AHL_2020__3__35_0/#DS}{Publisher reference}.
\bibitem[DY08]{DYaman} Dahmani, François; Yaman, Asli Symbolic dynamics and relatively hyperbolic groups, Groups Geom. Dyn., Volume 2 (2008) no. 2, pp. 165-184 \href{https://ahl.centre-mersenne.org/item/AHL_2020__3__35_0/#DYaman}{Publisher reference}.
\bibitem[Ger12]{Gera} Gerasimov, Victor Floyd maps for relatively hyperbolic groups, Geom. Funct. Anal., Volume 22 (2012) no. 5, pp. 1361-1399 \href{https://ahl.centre-mersenne.org/item/AHL_2020__3__35_0/#Gera}{Publisher reference}.
\bibitem[GRT]{GRT} Guentner, Erik; Reckwerdt, Eric; Tessera, R. Proper actions and weak amenability for classical relatively hyperbolic groups (in preparation) \href{https://ahl.centre-mersenne.org/item/AHL_2020__3__35_0/#GRT}{Publisher reference}.
\bibitem[GS18]{GS} Gruber, Dominik; Sisto, Alessandro Infinitely presented graphical small cancellation groups are acylindrically hyperbolic, Ann. Inst. Fourier, Volume 68 (2018) no. 6, pp. 2501-2552 \href{https://ahl.centre-mersenne.org/item/AHL_2020__3__35_0/#GS}{Publisher reference}.
\bibitem[Kaz67]{Kaz} Kazhdan, David A. On the connection of the dual space of a group with the structure of its closed subgroups, Funct. Anal. Appl., Volume 1 (1967) no. 1, pp. 63-65 \href{https://ahl.centre-mersenne.org/item/AHL_2020__3__35_0/#Kaz}{Publisher reference}.
\bibitem[Laf08]{La} Lafforgue, Vincent Un renforcement de la propriété (T), Duke Math. J., Volume 143 (2008) no. 3, pp. 559-602 \href{https://ahl.centre-mersenne.org/item/AHL_2020__3__35_0/#La}{Publisher reference}.
\bibitem[LS77]{LS77} Lyndon, Roger C.; Schupp, Paul E. Combinatorial group theory, Ergebnisse der Mathematik und ihrer Grenzgebiete, 89, Springer, 1977 \href{https://ahl.centre-mersenne.org/item/AHL_2020__3__35_0/#LS77}{Publisher reference}.
\bibitem[Min01]{Min} Mineyev, Igor Straightening and bounded cohomology of hyperbolic groups, Geom. Funct. Anal., Volume 11 (2001) no. 4, pp. 807-839 \href{https://ahl.centre-mersenne.org/item/AHL_2020__3__35_0/#Min}{Publisher reference}.
\bibitem[MS17]{MS} Martin, Alexandre; Steenbock, Markus A combination theorem for cubulation in small cancellation theory over free products, Ann. Inst. Fourier, Volume 67 (2017) no. 4, pp. 1613-1670 \href{https://ahl.centre-mersenne.org/item/AHL_2020__3__35_0/#MS}{Publisher reference}.
\bibitem[Nic13]{Ni} Nica, Bogdan Proper isometric actions of hyperbolic groups on L p -spaces, Compos. Math., Volume 149 (2013) no. 5, pp. 773-792 \href{https://ahl.centre-mersenne.org/item/AHL_2020__3__35_0/#Ni}{Publisher reference}.
\bibitem[NR97]{NiRe} Niblo, Graham; Reeves, Lawrence Groups acting on CAT(0) cube complexes, Geom. Topol., Volume 1 (1997), pp. 1-7 \href{https://ahl.centre-mersenne.org/item/AHL_2020__3__35_0/#NiRe}{Publisher reference}.
\bibitem[Pan90]{Pan} Pansu, Pierre Cohomologie L p  des variétés à courbure négative, cas du degré 1, Conference on partial differential equations and geometry (Rendiconti del Seminario Matematico), Universitá e Politecnico Torino, 1990, pp. 95-120 \href{https://ahl.centre-mersenne.org/item/AHL_2020__3__35_0/#Pan}{Publisher reference}.
\bibitem[Pan99]{Pankratev99} Pankratʼev, Anton E. Hyperbolic products of groups, Vestn. Mosk. Univ., Volume 1999 (1999) no. 2, pp. 9-13 \href{https://ahl.centre-mersenne.org/item/AHL_2020__3__35_0/#Pankratev99}{Publisher reference}.
\bibitem[Pul07]{Puls} Puls, Michael The first L p -cohomology of some groups with one end, Arch. Math., Volume 88 (2007) no. 6, pp. 500-506 \href{https://ahl.centre-mersenne.org/item/AHL_2020__3__35_0/#Puls}{Publisher reference}.
\bibitem[RS95]{RS95} Rips, Eliyahu; Sela, Zlil Canonical representatives and equations in hyperbolic groups, Invent. Math., Volume 120 (1995) no. 3, pp. 489-512 \href{https://ahl.centre-mersenne.org/item/AHL_2020__3__35_0/#RS95}{Publisher reference}.
\bibitem[Wis04]{Wise} Wise, Daniel Cubulating small cancellation groups, Geom. Funct. Anal., Volume 14 (2004) no. 1, pp. 150-214 \href{https://ahl.centre-mersenne.org/item/AHL_2020__3__35_0/#Wise}{Publisher reference}.
\bibitem[Yu05]{Yu} Yu, Guoliang Hyperbolic groups admit proper affine isometric actions on \(\ell\) p -spaces, Geom. Funct. Anal., Volume 15 (2005) no. 5, pp. 1144-1151 \href{https://ahl.centre-mersenne.org/item/AHL_2020__3__35_0/#Yu}{Publisher reference}.
\end{thebibliography}
\end{document}
